\section{Relation to prior work}
\label{sec:related}

\paragraph{Normalized persistence.}
Lowe et al.~\cite{lowe2026normalized} define normalized harmonic persistence
(their Problem~9, our \cref{def:problem}), prove that a low-energy relaxation
of it is \(\mathsf{DQC}_1\)-hard and in \(\BQP\), prove \(\SDQCone\)-hardness
of the local-Hamiltonian analogue, and conjecture \(\SDQCone\)-hardness of
Problem~9 (their Conjecture~1).  Their Conjecture~2 asks for a polynomial-time
map from history Hamiltonians \(H_1\) and \(H_2=H_1+\lambda H_{\rm out}\) to
weighted clique complexes \(X_1\subseteq X_2\) with two properties:
kernel fidelity and gap, \(\beta_d(X_i)=\dim\ker H_i\) together with
computable inverse-polynomial gap bounds; and functoriality,
\(\beta_d(X_1\to X_2)=\dim\ran(P_{\ker H_2}|_{\ker H_1})\).  Their Lemma~11
shows that Conjecture~2 implies Conjecture~1.  \Cref{thm:transfer,thm:quotient}
establish both properties for the exact \(G_2\) history Hamiltonians of
\cref{sec:circuit}, with \(H_{\rm out}\) at unit coefficient rather than at a
small coefficient \(\lambda\), with the induced map computed through the
quotient model, and with the harmonic encodings shown a posteriori to be
\(\eta/2\)-approximate isometries compatible with the inclusion
(\cref{eq:projector-closeness}).  \Cref{cor:sdqc} then proves Conjecture~1
over \(G_2\) in the fraction form.  The distinction between the two forms is
this: Definition~2 of Lowe et al.\ imposes the decision condition on the
acceptance probability \(p_x\) of the maximally mixed state, whereas the
quantity estimated by the oracle in the proofs of their Theorems~5 and~7 is
the fraction \(f_x=\dim S_x/2^{q-1}\), and after the Marriott--Watrous
amplification used there the acceptance probability of the amplified circuit
is within \(1/\poly\) of \(f_x\) (their equation~(9)) but no longer within
\(1/\poly\) of \(p_x\).  Since \(f_x\le p_x\le(1+2f_x)/3\)
(\cref{app:source-boundaries}), the trace promise implies the fraction
promise exactly when \(b<(3a-1)/2\), and \cref{cor:sdqc} covers those
thresholds in the trace form as well; for other thresholds no reduction whose
target depends only on \(S_x\) can exist.  Their containment result
(Lemma~7) needs a preparation circuit for the uniform mixture over
\(\ker\Delta_{1,d}\) and an inverse-polynomial lower bound on the nonzero
eigenvalues of \(P_{\cH_2}P_{\cH_1}P_{\cH_2}\); \cref{sec:advantage} proves
both on our instances, which is why \cref{thm:advantage} can be stated in
the same form as their Theorems~5, 7 and~8.

\paragraph{Clique homology gadgets.}
Crichigno and Kohler~\cite{crichigno2024clique} give a parsimonious reduction
in which satisfying states correspond one-to-one to holes, so exact
multiplicity by itself is not new; their instances carry no gap promise.
King and Kohler~\cite{king2026gapped} introduce the weighted gadgets and the
many-gadget spectral analysis behind gapped clique homology, proving a gap
when the logical Hamiltonian is positive definite, with a gadget weight
proportional to the logical gap; Hayakawa~\cite{hayakawa2026unweighted}
removes the weights under the same positive-definite promise.  Our transfer
theorem is formulated from finite zero-weight certificates, applies to
arbitrary geometric chains and to degenerate kernels, decouples
\(\lambda\) from \(g\), and takes its exponent from the fillings rather than
from the locality.  Its padding step differs from the coordinate-bulk
formulation displayed in the arXiv version of King and
Kohler~\cite[Claim~10.2]{king2023gappedv2}: there the bulk consists of all
coordinates containing a private vertex, which in the padded complex also
contains boundary terms obtained by deleting a weight-one outside vertex,
whereas we keep local private coordinates tensored with outside harmonics
in every bidegree (\cref{lem:scaling-padding}).  We also work with the fixed
local matrices before tensoring with the outside identity, so that no
operator bound on a projector of growing rank has to be inferred from
basiswise perturbation estimates.  These are differences of proof technique,
and we make no claim about the presentation in the published
version~\cite{king2026gapped}.

\paragraph{Counting Betti numbers.}
Crichigno and Kohler prove that computing Betti numbers of clique complexes
is \(\sharpP\)-hard~\cite[Theorem~4]{crichigno2024clique}, through the
\(\sharpP\)-hardness of counting quantum \(4\)-SAT solutions; their instances
have no gap promise.  Schmidhuber and Lloyd prove \(\sharpP\)-hardness of
exact Betti numbers and \(\mathsf{NP}\)-hardness of multiplicative
approximation for clique complexes, even in the clique-dense regime where
the quantum algorithms perform best~\cite{schmidhuber2023limitations}; the
restriction is a density condition, not a spectral one.  Cade and Crichigno
prove \(\sharpP\)-completeness with an inverse-polynomial gap for general chain
complexes with \(6\)-local coboundaries~\cite[Theorem~5]{cade2024complexity},
using the counting class \(\#\BQP\) of Brown, Flammia, and
Schuch~\cite{brown2011computational}, whose definition includes a gap promise
on the acceptance operator and which equals \(\sharpP\) under weakly
parsimonious reductions.  \Cref{thm:gapped-betti} is the clique-complex
statement with the gap, for weighted and unweighted graphs; it is a direct
corollary of the degenerate-kernel transfer, since the classical source has
\(M_\varphi|_{S^\perp}=0\).

\paragraph{Persistence hardness.}
Gyurik et al.~\cite{gyurik2026provable} prove \(\BQPone\)-hardness, over the
gate set \(\{\mathrm{CNOT},U_{\rm Pyth}\}\), of deciding whether a
\emph{given} harmonic representative of \(X_1\) survives in \(X_2\), with a
gap promise on the final Laplacian.  Normalized persistence asks instead for
the fraction of the full initial homology that survives, with no
representative given, so a hardness proof must control the complete kernels
at both endpoints.  Persistent Laplacians~\cite{memoli2022persistent} give a
spectral formulation of persistent Betti numbers; we use ordinary endpoint
Laplacians for the promise and the quotient model for the map.

\paragraph{Sources and gate sets.}
Rudolph~\cite{rudolph2026universal} defines the gate-dependent classes,
proves exact error reduction inside gate sets containing \(G_2\), shows
\(\BQPone^{G}\subseteq\BQPone^{G_2}\) for every finite \(G\) with entries in
some \(\mathbb Q(\zeta_{2^k})\), and supplies the gain/loss split of
\(H\otimes H\) with rational three-term constraint vectors and the base
clique gadgets used in our palette.  Cade and
Crichigno~\cite{cade2024complexity} relate cohomology problems to
\(\QMAone\) and \(\mathsf{DQC}_1\), and Gyurik, Cade, and
Dunjko~\cite{gyurik2022towards} connected quantum topological data analysis
to \(\mathsf{DQC}_1\)~\cite{knill1998power}.  The history-state construction
with a unary clock and neighbor guards is
Kitaev's~\cite{kitaev2002classical,kempe2006complexity}, with the rational
amplitude bookkeeping of \cref{sec:circuit}; the amplification used by Lowe
et al.\ is that of Marriott and Watrous~\cite{marriott2005quantum}.

\paragraph{Unweighting.}
Hayakawa's symmetric/asymmetric decomposition of the common-copy
blow-up~\cite[Lemma~4.2 and Theorem~4.3]{hayakawa2026unweighted} is used as a
black box.  Our only addition is the observation that identical labeled copy
blocks at both filtration levels make the symmetric isometries commute with
inclusion, \cref{eq:unweight-natural}, which also transfers the containment
conditions.  We do not count this as a separate contribution.
